\begin{figure}[t]
\centering
\begin{tikzpicture}
\begin{axis}[
  width=0.60\textwidth, height=4.2cm,
  xmin=0.7, xmax=6.3, ymin=0, ymax=105,
  xtick={1,2,3,4,5,6},
  xticklabels={session note,consolidated,resum 130,resum 110,resum 90,resum 70},
  x tick label style={rotate=28,anchor=east,font=\scriptsize},
  ylabel={chains still stating the qualifier (\%)}, ylabel style={font=\scriptsize},
  tick label style={font=\scriptsize},
  ymajorgrids, grid style={gray!20}, axis lines*=left,
  legend style={font=\scriptsize,at={(0.03,0.06)},anchor=south west,draw=none,fill=none},
]
\addplot[mark=*,mark size=1.8pt,blue!70!black,thick] coordinates {\FiveErosionCoords};
\addlegendentry{target memory (qualifier stated)}
\addplot[mark=square*,mark size=1.8pt,orange!85!black,thick,dashed] coordinates {\FiveScopeErosionCoords};
\addlegendentry{scope-restricted memories (qualifier stated)}
\end{axis}
\end{tikzpicture}
\caption{What six generations of benign consolidation removed, judged by a
frozen model scorer (\FiveScorerModel, itself one of the six consolidator
models) with \FiveFalsePosAll{} false positives on records carrying no
qualifier. The corruption our controlled sections install --- a body with no
negative content left --- never occurs (\FivePurelyPositive{} chain-generations).
Plotted is the scorer's judgment that the body still \emph{states} its
qualifier; the quantified negative itself never falls (\FiveNumAllGens{} at
every generation). What erodes monotonically \emph{within} the tested horizon
is the qualifier as stated: the target memory falls to
\FiveTgtGenSixK/\FiveTgtGenSixN{} and scope restrictions to
\FiveScopeGenSixK/\FiveScopeGenSixN. This is a decline observed inside the
tested range, not a plateau, and it is why we do not extrapolate the negative
result to longer horizons.}
\label{fig:erosion}
\end{figure}
